\documentclass[../Main.tex]{subfiles}

\begin{document}

\section{Reflexion} \label{sec:Reflexion}

\subsection{Reflexion des Arbeitsprozesses}

Dieses Projekt hat etwa ein Jahr unserer Zeit in Anspruch genommen und mehrere Phasen durchlaufen, die wir im Folgenden der Reihe nach aufgreifen.

Ganz am Anfang waren wir noch völlig unsicher und hatten ganz andere Vorstellungen davon, welches Produkt am Ende entstehen würde. Also stocherten wir in alle Bereiche hinein: Wir begannen mit Spielsystemen, von denen wir ohne Recherche annahmen, dass wir sie brauchen, überlegten, wie wir die Erkenntnisfrage beantworten wollen, und recherchierten, wie man recherchiert. Bis zur Zwischenpräsentation blieben uns noch vier Monate. Wir waren wenig zielorientiert und machten einfach, was uns gefiel oder was wir noch nicht ausprobiert hatten. Das mag planlos wirken, rückblickend war es aber ein Auskundschaften dessen, was es gibt. Die Fallhöhe war zudem gering: Was auch immer wir machten, nichts musste verworfen werden, und die Arbeitsschritte hingen noch kaum voneinander ab. Dass das aufging, war zwar eher ein glücklicher Zufall als strukturiertes Denken unsererseits, doch so ist die Natur eines Projektanfangs. Unter Leuten, die Projekte machen, heisst es ja, sie hätten nie angefangen, hätten sie gewusst, wie viel Arbeit es wird. Könnten wir noch einmal beginnen, würden wir es genauso angehen, mit dem einzigen Unterschied, dass wir wüssten, dass wir am Anfang eine Narrenfrist haben.

Bis zur Zwischenpräsentation spielte die Erkenntnisfrage eine untergeordnete Rolle, denn die dafür relevanten Systeme beruhten auf einer technischen Basis, die wir noch nicht hatten: einem Spielercharakter mit Animationen und Steuerung, Gegnern und Code, der das Spiel überhaupt zum Starten bringt. Obwohl wir dabei die Recherche vernachlässigten, war das rückblickend die richtige Entscheidung. Zum einen machte uns dieser Teil am meisten Freude, zum anderen mussten wir uns zuerst eine Vorstellung von unserem eigenen Spiel bilden. Erst nachdem wir die ersten Systeme gebaut hatten, wussten wir, worauf wir in anderen Spielen achten mussten, und konnten erst recht eine gute Recherche machen. Bei einem nächsten Projekt würden wir die Reihenfolge davon abhängig machen, wie konkret unsere Vorstellungen sind und wie gross unser Kenntnisstand ist: Haben wir beides, recherchieren wir zuerst. Fehlt eines davon, sammeln wir zuerst eigene Erfahrungen -- also genau das, was wir diesmal eher unbeabsichtigt getan haben.

Gleichzeitig rückte die Zwischenpräsentation näher. Features, die wir zwar wollten, die für die Erkenntnisfrage aber wenig relevant waren, mussten wir schon jetzt streichen oder stark eingrenzen, und wir mussten priorisieren, was für ein Spiel essenziell ist. Dazu kam ein begründet geplanter Playtest, und die Checkliste aus dem Vertrag hatte noch viele offene Punkte. Wir haben alles erreicht, was wir vertraglich festgehalten hatten, auch wenn es sehr, sehr knapp wurde. Nicht das Fertigstellen, sondern die Qualitätssicherung hat uns einen Strich durch die Rechnung gemacht, weil wir möglichst wenige Kompromisse eingehen wollten. Im Nachhinein wäre es interessant gewesen zu erfahren, wie 80 oder 85 statt 100 Prozent unseres Anspruchs bewertet worden wären. Da die Zwischenbewertung nicht in die Note der Maturitätsarbeit einfliesst, ist das eine verpasste Chance. Andererseits konnten wir Erfahrung mit Crunchen sammeln, die uns am Ende ermöglicht hat, mit genügend Zeit und kaum Crunchen 95 bis 100 Prozent herauszuholen.

Zur Zwischenpräsentation erhielten wir eine Rückmeldung: Vieles war gut, einiges weniger. Die Kritik betraf weniger den Prototyp als alles rundherum, besonders die begleitenden Texte. Beim Umsetzen hatten wir zunächst Schwierigkeiten, denn hätten wir die bemängelten Punkte verinnerlicht gehabt, hätten wir sie von Anfang an richtig gemacht. Deshalb stellten wir unserer Betreuungsperson so lange Fragen, bis wir sicher waren, sie verstanden zu haben.

Nach der Zwischenpräsentation hatten wir wieder genug Zeit, und die Abgabe war weit genug entfernt, dass wir unsere Ambitionen beim Projektumfang vergrösserten. Gleichzeitig nahm das anfängliche Tempo ab, mit dem wir Neues produzierten. Zum einen war das Projekt inzwischen so gross, dass kleinere Neuerungen wenig Unterschied machten. Zum anderen war unser anfänglicher Eifer zwar weit davon entfernt, erloschen zu sein, doch das Projekt zu pflegen und Aufgaben zu erledigen, die uns weniger Spass machten und sich mittlerweile angehäuft hatten, setzte uns langsam zu.

Mittlerweile hingen die einzelnen Arbeitsschritte so stark voneinander ab, dass wir in unserer Partnerarbeit nicht mehr einfach tun und lassen konnten, was wir wollten, sondern uns aufeinander abstimmen mussten. Wir arbeiteten deshalb an möglichst grossen, voneinander unabhängigen Blöcken und vereinbarten ein Datum, wenn sich unsere Arbeiten kreuzten. In dieser Zeit entstand auch fast die gesamte restliche Recherche. Das funktionierte sehr gut, wenn beide rechtzeitig fertig wurden, was allerdings nicht immer der Fall war. Künftig würden wir deshalb einen dritten, grösseren Arbeitsstrang offenhalten, an dem weiterarbeiten kann, wer gerade durch die andere Person blockiert ist.

Nach Abschluss der Recherche wussten wir genau, welche Systeme noch implementiert werden mussten, und mit der näher rückenden Abgabe hatten wir nun auch die Beantwortung der Erkenntnisfrage im Blick. Hier beginnt der wirklich gut geplante Teil unserer Arbeit. Da wir nun ein Fundament hatten, das in eine klare Richtung ging, und wussten, dass unsere Arbeit nicht mehr verschwendet sein würde, konnten wir zielgerichtet auf die Erkenntnisfrage hinarbeiten. Dieser Teil ist uns besonders gelungen, weil wir auf Hinweis unserer Betreuungsperson einen Plan bis zur Abgabe erstellt haben. Dass dieser Plan gut wurde, verdanken wir den Erfahrungen aus der Zwischenpräsentation, bei der wir bereits gelernt hatten, wie man auf eine Abgabe hinarbeitet und plant.

Unsere ersten ersten fertigen export hatten wir am 25. dann vertig spielete sich aber grausam, obwohl techinisch alle elemente der recherche umgesetzt wurden. es war allerdings halb so schlimm wie es zu anfagns den schein gehabt hatte, da wir nur ein paar stellstrauben nachjustieren musten und das spiel war unsere recherche gerecht und bereit für ten playtest. diese erfahrung hat uns gezeigt, dass der teufel im detail liegt und das problem ist genau entstanden, als recherche und protoyp komplett ineinander hinein gegriffen haben. sie haben ahlt einfach nicht so gut gegriffen. aber mit etwas nachjustieren ging es. wir haben bis hier die recherche und den prototyp seperat angesehen, ws dazu geführt hat. anfans war dieser gedanke nicht ganz falsch aber je weiter wir gearbeitet haben, desto mehr hatten die beiden miteinenander zu tun. beim nächsten mal wissen wir, dass wir darauf achten müssen wenn solche zwei teile zusammen kommen, dass siw wirklich gut aufeinanger schon im vorhinein abgestimmt werden sollen und nicht als blosser gedanke danach.

Gegen Ende der Arbeit führten wir die Playtests durch. Sie verliefen recht gut, da wir für die Recherche bereits selbst Playtests durchgeführt hatten. In der Schlussphase bereiteten uns eher Formalitäten Sorgen. Auch diese Phase verlief unserer Einschätzung nach ziemlich gut, nicht weil wir darin etwas Besonderes gemacht hätten, sondern weil wir zuvor gut geplant und uns genug Zeit dafür gelassen hatten. Zudem war uns bewusst, dass es knapp wird, wenn wir uns nicht an den Plan halten, auch wenn es noch nach genug Zeit aussah.

\begin{itemize}[leftmargin=1.5cm]
    \item \textcolor{red}{\textbf{Wendepunkt fehlt} (Journal 24./25.10., 13~h): Obwohl der Prototyp technisch fertig war, machten die Kämpfe bei euren eigenen Tests wenig Spass. Ihr habt die Encounter gekürzt, die Gegner aggressiver gemacht und noch nicht umgesetzte Rechercheergebnisse nachgeholt. Eure Erkenntnis im Journal: \enquote{Die technische Fertigstellung allein garantiert noch kein gutes Spielerlebnis}, und weil ihr das Level so oft gespielt hattet, brauchte es externe Playtests. Gehört vor den Absatz zu den Playtests; dazu, was ihr anders machen würdet (z.\,B. früher extern testen).}
    \item \textcolor{red}{\textbf{Plan und Vertrag:} Laut Vertrag waren die Playtests für Mitte Oktober und die Verbesserungen für Ende Oktober geplant; tatsächlich fanden die Playtests vom 26.10.\ bis 2.11.\ statt, und die Verbesserungen stehen nur im Bericht. Das passt nicht ganz zu \enquote{wirklich gut geplant} -- kurz erklären, warum es sich verschoben hat und warum ihr die Planung trotzdem als gelungen seht.}
    \item \textcolor{red}{\textbf{Scope konkret machen:} 1--2 Beispiele für Features, die ihr vor der Zwischenpräsentation gestrichen habt (Vertrag bereinigt am 20.5.2026), und was nach der Zwischenpräsentation dazugekommen ist. Dazu ein Urteil: War die Vergrösserung des Projektumfangs im Rückblick klug?}
    \item \textcolor{red}{\textbf{Rückmeldung der Zwischenbewertung:} 1--2 konkrete Punkte nennen, die ihr danach geändert habt (die Bewertung verlangte u.\,a.\ bessere Belege, vollständige KI-Dokumentation, vertiefte Reflexionen, Zeitaufwand dokumentieren).}
    \item \textcolor{red}{\textbf{Motivation:} Wie seid ihr damit umgegangen, dass euch die angehäuften Aufgaben zusetzten?}
    \item \textcolor{red}{\textbf{Playtests präzisieren:} Welche eigenen Playtests \enquote{für die Recherche} meint ihr -- den Playtest vor der Zwischenpräsentation oder eure eigenen Tests am Level?}
    \item \textcolor{red}{\textbf{Verweis aufs Arbeitsjournal}, wo der Verlauf im Detail dokumentiert ist.}
\end{itemize}

\subsection{Gruppenarbeit und Aufgabenteilung}

\begin{itemize}[leftmargin=1.5cm]
    \item \textcolor{red}{\textbf{1. Wie ihr die Arbeit aufgeteilt habt.} Laut Zwischenbewertung: Rodin stärker bei Animation, UI, Bericht und README; Zhile stärker bei Architektur, Bewegung, State Machines und Kampf-Grundlagen. Das deckt sich mit Git: Zhile vor allem \code{Entities/Player}, \code{Common} (Utility, Tools), \code{Core}, Kamera; Rodin vor allem \code{Interfaces} und Menüs sowie die TrenchBroom-Karten des Levels.}
    \item \textcolor{red}{\textbf{2. Tabelle Aufgabenteilung und Eigenleistung} für Prototyp und Bericht (Rodin / Zhile / gemeinsam). Zum Bericht: Kapitel 1--8 und die \game{NieR}-Analyse Rodin, die \game{Ninja-Gaiden}-Analyse Zhile, README und Architekturdiagramm (Juni) vor allem Zhile. Bitte prüfen und ergänzen (Playtests, Leveldesign, Audio, \ldots).}
    \item \textcolor{red}{\textbf{3. Zusammenarbeit und Kommunikation:} Wie habt ihr euch abgesprochen (Treffen, Discord, GitHub)? Wie habt ihr Entscheidungen getroffen?}
    \item \textcolor{red}{\textbf{4. Veränderung im Verlauf:} z.\,B. ab Oktober Rodin vor allem Level, Playtests und Bericht, Zhile weiter am Spiel (laut Git im Oktober sehr aktiv). Wie kam es dazu?}
    \item \textcolor{red}{\textbf{5. Fazit:} Was lief gut, was war schwierig, was würdet ihr in einer Gruppenarbeit anders machen? Hinweis: Commit-Zahlen sagen wenig über den Aufwand aus, deshalb nicht als Massstab verwenden.}
\end{itemize}

\subsection{Verwendung von KI} \label{sec:ReflexionKI}

\begin{itemize}[leftmargin=1.5cm]
    \item \textcolor{red}{\textbf{1. Welche Tools wofür:} Code (Treppen-Werkzeug und Animations-Hilfsskript im Frühling, Audiosystem im Oktober), Begleitbericht (Korrektur, Struktur, Rückmeldungen). Details in der KI-Deklaration (Abschnitt~\ref{sec:KIDeklaration}).}
    \item \textcolor{red}{\textbf{2. Was geholfen hat:} z.\,B. Widersprüche zwischen Kapiteln finden, unbelegte Aussagen markieren, \LaTeX-Formatierung, Fehlersuche im Code.}
    \item \textcolor{red}{\textbf{3. Grenzen und Probleme:} Die KI ging mehrmals von falschen Annahmen aus, die ihr korrigieren musstet, z.\,B. dass Themenfindung mit KI gemacht wurde, dass ein Tutorial geplant war, welche Räume Brückenräume sind, oder dass in den Protokollen ein anderer Archer gemeint war.}
    \item \textcolor{red}{\textbf{4. Wie ihr die Ergebnisse überprüft habt:} Abgleich mit Quellen, Playtest-Protokollen und dem Spiel selbst; Vorschläge nicht ungeprüft übernommen, einige bewusst verworfen.}
    \item \textcolor{red}{\textbf{5. Eigenanteil sichern:} Interpretation und Urteile selbst festgelegt; von der KI formulierte Stellen markiert und in eigenen Worten neu geschrieben; Regel \enquote{nicht wortwörtlich übernehmen} aus dem Merkblatt.}
    \item \textcolor{red}{\textbf{6. Fazit:} Würdet ihr KI wieder so einsetzen? Was würdet ihr anders machen, z.\,B. KI-Nutzung von Anfang an laufend im Journal dokumentieren?}
\end{itemize}

\subsection{Persönliche Erkenntnisse}

\subsubsection{Rodin Arslan}

\begin{itemize}[leftmargin=1.5cm]
    \item \textcolor{red}{\textbf{1. Fachlich:} z.\,B. Animation (u.\,a. Hängeanimation), UI und Menüs, Leveldesign mit TrenchBroom, Gegner- und Encounter-Design, Playtesting, \LaTeX\ und BibTeX.}
    \item \textcolor{red}{\textbf{2. Methodisch:} Spielanalyse mit dem MDA-Framework, Recherche und Zitieren, Playtests planen und auswerten, Projektplanung.}
    \item \textcolor{red}{\textbf{3. Persönlich:} z.\,B. dass ein technisch fertiges Spiel noch nicht gut spielbar ist (Journal 24./25.10.), Zeitmanagement in der Schlussphase, Umgang mit Rückmeldungen.}
    \item \textcolor{red}{\textbf{4. Was du mitnimmst} für die Zukunft.}
\end{itemize}

\subsubsection{Zhile Zhu}

\textcolor{red}{ZHILE DAS IST FÜR DICH}

\end{document}
